\documentclass[a4paper]{article}

\usepackage[utf8]{inputenc}
\usepackage{fullpage}
\usepackage[T1]{fontenc}
\usepackage[serbian]{babel}
\usepackage{graphicx}
\usepackage{amsmath}
\usepackage{booktabs}
\usepackage{url}
\usepackage{hyperref}
\usepackage{xcolor}
\usepackage{float}
\usepackage{caption}
\usepackage{subcaption}
\usepackage{csquotes}
\usepackage{tabularx}
\usepackage{listings}
\hypersetup{colorlinks=true, linkcolor=blue, urlcolor=cyan, citecolor=blue}

\definecolor{codebg}{rgb}{0.97,0.97,0.97}
\definecolor{codecomment}{rgb}{0.4,0.4,0.4}
\definecolor{codekeyword}{rgb}{0.0,0.3,0.7}
\definecolor{codestring}{rgb}{0.2,0.5,0.2}

\lstset{
    backgroundcolor=\color{codebg},
    basicstyle=\ttfamily\small,
    commentstyle=\color{codecomment},
    keywordstyle=\color{codekeyword}\bfseries,
    stringstyle=\color{codestring},
    breaklines=true,
    frame=single,
    rulecolor=\color{gray!30},
    showstringspaces=false,
    numbers=left,
    numberstyle=\tiny\color{gray},
    numbersep=6pt,
    xleftmargin=1.5em,
    xrightmargin=0.5em,
}

\lstdefinelanguage{Kotlin}{
    keywords={fun, val, var, if, else, while, for, return, import, class, object,
              interface, when, is, in, null, true, false, do, break, continue, throw},
    morecomment=[l]{//},
    morecomment=[s]{/*}{*/},
    morestring=[b]",
    sensitive=true
}

\lstdefinelanguage{CSharp}{
    keywords={public, private, protected, internal, sealed, partial, override,
              class, interface, using, namespace, return, if, else, throw, null,
              var, void, string, int, bool, new, this, base},
    morecomment=[l]{//},
    morecomment=[s]{/*}{*/},
    morestring=[b]",
    sensitive=true
}
\renewcommand{\lstlistingname}{Isečak koda}

\begin{document}

\title{Podrška za Kotlin u okviru LINVAST biblioteke\\
\small{Opis sistema u okviru seminarskog rada\\
Matematički fakultet, Univerzitet u Beogradu}}
\author{Anđela Jovanović\\
\texttt{mi251035@alas.matf.bg.ac.rs}}
\date{Jun 2026.}

\maketitle
\begin{abstract}
Ovaj rad opisuje proširenje biblioteke LINVAST podrškom za programski jezik
Kotlin. LINVAST je .NET biblioteka koja definiše jezički invarijantno
apstraktno sintaksičko stablo (AST) za imperativne jezike, omogućavajući
analizu i poređenje programa pisanih u različitim jezicima kroz jedinstvenu
reprezentaciju.
\end{abstract}
\tableofcontents
\newpage

% -----------------------------------------------------------------------
\section{Opis problema}
% -----------------------------------------------------------------------

U ovom poglavlju se opisuju motivacija za uvođenje jezički invarijantne statičke
analize i konkretni ciljevi koji se ovim radom ostvaruju.

\subsection{Motivacija}

Alati za statičku analizu programa tipično su namenjeni jednom konkretnom
programskom jeziku. Kada je potrebno primeniti istu vrstu analize, kao što su
detekcija mrtvog koda, izračunavanje konstantnih izraza ili provera ispravnosti
kontrolnog toka, na programe pisane u različitim jezicima, razvojni timovi su
prinuđeni da održavaju zasebne, jezički specifične alate, ili da značajno suzuju
obuhvat same analize.

\vspace{2px}
Apstraktno sintaksičko stablo (skr.~AST) nastaje kao rezultat sintaksičke
analize ulaznog programa i čini osnovu za semantičku analizu i kasniju fazu
kompilacije~\cite{ristovic2020}. Budući da je AST vezan za gramatiku konkretnog
programskog jezika, programi pisani u različitim jezicima ne mogu se direktno
porediti posmatranjem njihovih stabala, što je posebno ograničavajuće u
kontekstu migracije na nove tehnologije ili provere ekvivalentnosti implementacija.

\vspace{2px}
\textbf{LINVAST} (Language-INVariant AST) je .NET biblioteka koja rešava
ovaj problem definisanjem \emph{jedinstvenog apstraktnog sintaksičkog stabla}
sposobnog da predstavi programe iz više imperativnih programskih jezika. Za svaki
podržani jezik postoji poseban \emph{graditelj} (engl.\ \emph{builder}) koji
parsira izvorni kod i preslikava ga u ovu zajedničku reprezentaciju. Jednom kada
se program prevede u LINVAST oblik, konstrukti poput deklaracija, izraza,
kontrolnog toka i funkcija mogu se upoređivati, transformisati ili analizirati
jezički nezavisnih posetiocima (engl.\ \emph{visitors}), bez ikakve svesti o
izvornom programskom jeziku~\cite{ristovic2020}.

\subsection{Cilj ovog rada}

Cilj ovog seminarskog rada je \textbf{proširiti LINVAST podrškom za Kotlin}.
Kotlin je moderan, statički tipiziran jezik koji se izvršava na JVM platformi;
njegova gramatika obuhvata konstrukte koji ne postoje u jezicima koje LINVAST
već podržava (C, Lua, Go, Java). Konkretno, proširenje zahteva:

\begin{enumerate}
    \item Preuzimanje ANTLR~4 gramatike za Kotlin i integrisanje generisanog
          C\# leksera i parsera u projekat.
    \item Implementaciju klase \texttt{KotlinASTBuilder}, raspoređene u šest
          parcijalnih \texttt{.cs} datoteka, koja obilazi generisano stablo
          parsiranja i proizvodi LINVAST AST čvorove.
    \item Verifikaciju ispravnosti kroz postojeću testnu infrastrukturu.
\end{enumerate}

Implementacija pokriva: deklaracije \texttt{val}/\texttt{var}
svojstava, deklaracije funkcija na nivou datoteke, direktive uvoza
(\texttt{import}), osnovni kontrolni tok (\texttt{if}/\texttt{else},
\texttt{while}, \texttt{for}), \texttt{return} naredbe, kompletnu hijerarhiju
aritmetičko-logičkih izraza, kao i deklaracije klasa i interfejsa sa primarnim
konstruktorom, naslednim tipovima i telom klase. 

Lambda izrazi, ekstenzivne funkcije, koroutine i napredni tipski konstrukti izvan su obuhvata
i razmatraju se kao budući rad u poglavlju~\ref{sec:future}.

\newpage
% -----------------------------------------------------------------------
\section{Arhitektura sistema}
% -----------------------------------------------------------------------

U ovom poglavlju opisuje se kako je LINVAST rešenje strukturisano i gde
se Kotlin builder uklapa u tu celinu.

\subsection{Struktura rešenja}

LINVAST rešenje sastoji se od četiri .NET projekta.

\begin{table}[H]
\centering
\begin{tabularx}{\textwidth}{lX}
\toprule
\textbf{Projekat} & \textbf{Sadržaj} \\
\midrule
\texttt{LINVAST} &
    Osnovna biblioteka: apstraktna klasa \texttt{ASTNode},
    interfejsi \texttt{IASTBuilder<TParser>} i \texttt{IAbstractASTBuilder},
    atribut \texttt{ASTBuilderAttribute} za automatsko otkrivanje,
    tipovi izuzetaka i podrška za beleženje. \\
\addlinespace
\texttt{LINVAST.Imperative} &
    Konkretni tipovi AST čvorova za imperativne jezike (\texttt{SourceNode},
    \texttt{FuncNode}, \texttt{DeclStatNode}, \texttt{IfStatNode}, \ldots) i
    svi jezički specifični graditelji: C, Lua, Go, Java, \textbf{Kotlin}. \\
\addlinespace
\texttt{LINVAST.Imperative.Visitors} &
    Višekratno upotrebljivi posetioci AST-a: \texttt{ExpressionEvaluator}
    (sažimanje konstantnih izraza), \texttt{ConstantExpressionEvaluator},
    \texttt{SymbolicExpressionBuilder}. \\
\addlinespace
\texttt{LINVAST.Tests} &
    Testna skripta zasnovana na NUnit~3, organizovana po direktorijumima
    za svaki jezik; svaka klasa testova nasleđuje zajedničku osnovnu klasu
    koja pruža pomoćne metode za proveravanja. \\
\bottomrule
\end{tabularx}
\caption{Struktura projekata u LINVAST rešenju}
\end{table}

\subsection{Modul Kotlin builder-a}

Kotlin proširenje nalazi se u direktorijumu
\texttt{LINVAST.Imperative/Builders/Kotlin/} i prati isti raspored koji koristi
svaki drugi jezički graditelj u projektu:

\begin{verbatim}
Builders/Kotlin/
+-- ANTLR/
|   +-- KotlinLexer.g4             (lekser gramatika)
|   +-- KotlinParser.g4            (parser gramatika)
|   +-- KotlinLexer.cs             (generisani lekser)
|   +-- KotlinParser.cs            (generisani parser)
|   +-- KotlinParserVisitor.cs     (generisani interfejs posetioca)
|   +-- KotlinParserBaseVisitor.cs (generisana osnovna klasa)
+-- KotlinASTBuilder.cs            (ulazna tacka, deklaracija klase)
+-- KotlinASTBuilder.Declarations.cs
+-- KotlinASTBuilder.Functions.cs
+-- KotlinASTBuilder.Statements.cs
+-- KotlinASTBuilder.Expressions.cs
+-- KotlinASTBuilder.Types.cs
\end{verbatim}

\begin{sloppypar}
\texttt{KotlinASTBuilder} je \texttt{sealed partial} klasa koja nasleđuje
ANTLR-generisanu klasu \texttt{KotlinParserBaseVisitor<ASTNode>} i implementira
interfejs \texttt{IASTBuilder<KotlinParser>}. Šest izvornih datoteka kompiliraju
se zajedno putem mehanizma C\# \emph{delimičnih klasa} (engl.\ \emph{partial
class}), pri čemu svaka pokriva zasebnu oblast odgovornosti.
\end{sloppypar}

\newpage
\subsection{Automatsko otkrivanje}

\texttt{ImperativeASTFactory} koristi refleksiju za pronalaženje graditelja u
vreme izvođenja. Svaka klasa anotirana sa \texttt{[ASTBuilder(".kt")]}
automatski se bira kada se datoteka sa ekstenzijom \texttt{.kt} prosledi metodi
\texttt{BuildFromFile}.

\subsection{Tok obrade}

Celokupan put od izvornog Kotlin koda do LINVAST stabla prolazi kroz pet
uzastopnih koraka prikazanih ispod:

\begin{center}
\texttt{.kt datoteka}
$\;\to\;$ \texttt{CharStream}
$\;\to\;$ \texttt{KotlinLexer}
$\;\to\;$ \texttt{CommonTokenStream}
$\;\to\;$ \texttt{KotlinParser}
$\;\to\;$ \texttt{KotlinASTBuilder}
$\;\to\;$ \texttt{ASTNode}
\end{center}

\vspace{5px}
Metoda \texttt{CreateParser} realizuje prvu polovinu ovog puta:

\begin{lstlisting}[language=CSharp, caption={Kreiranje parsera iz izvornog koda}]
public KotlinParser CreateParser(string code)
{
    ICharStream stream  = CharStreams.fromstring(code);
    var lexer           = new KotlinLexer(stream);
    lexer.AddErrorListener(new ThrowExceptionErrorListener());

    ITokenStream tokens = new CommonTokenStream(lexer);
    var parser          = new KotlinParser(tokens);
    parser.RemoveErrorListeners();
    parser.AddErrorListener(new ThrowExceptionErrorListener());

    return parser;
}
\end{lstlisting}

\vspace{5px}
\noindent
Ulazni tekst se najpre obmotava u \texttt{CharStream} kako bi ANTLR mogao da
ga čita karakter po karakter. \texttt{KotlinLexer} zatim vrši tokenizaciju, to jest
grupiše karaktere u tokene poput ključnih reči (\texttt{val}, \texttt{fun}),
operatora (\texttt{+}, \texttt{<}), literala i identifikatora.
\texttt{CommonTokenStream} baferuje token niz kako bi parser mogao gledati
unapred. \texttt{KotlinParser} čita niz tokena i gradi konkretno sintaksičko
stablo (CST) u skladu sa pravilima Kotlin gramatike. Prilagođeni
\texttt{ThrowExceptionErrorListener} na oba mesta zamenjuje ANTLR-ov
podrazumevani listener koji štampa greške na konzoli - svaka sintaksička
ili leksička greška se odmah pretvara u izuzetak.

\vspace{5px}
Drugu polovinu puta realizuje metoda \texttt{BuildFromSource}:
\begin{lstlisting}[language=CSharp, caption={Izgradnja LINVAST stabla iz izvornog koda}]
public ASTNode BuildFromSource(string code)
{
    return this.Visit(this.CreateParser(code).kotlinFile());
}
\end{lstlisting}

\vspace{5px}
\noindent
Poziv \texttt{.kotlinFile()} nalaže parseru da krene od korenskog gramatičkog
pravila koje pokriva čitavu datoteku. \texttt{this.Visit(\ldots)} zatim
pokreće obilazak dobijenog CST-a pomoću metoda posetioca definisanih u
\texttt{KotlinASTBuilder}-u i njegovim delimičnim datotekama, i vraća koren
LINVAST stabla.

\newpage
% -----------------------------------------------------------------------
\section{Opis rešenja}
% -----------------------------------------------------------------------

U ovom poglavlju detaljno se opisuju pristup implementaciji, preslikavanje
pravila Kotlin gramatike u LINVAST čvorove i ključne odluke donete tokom razvoja.

\subsection{Obrazac Posetilac i obilazak CST-a}
ANTLR~4 generiše klasu \texttt{KotlinParserBaseVisitor<T>} koja implementira
\textbf{obrazac Posetilac} (engl.\ \emph{Visitor pattern}) nad konkretnim
sintaksičkim stablom: za svako pravilo gramatike postoji odgovarajuća
\texttt{Visit*} metoda. Svaka metoda koja nije redefinisana baca izuzetak
čime se garantuje da svaki neimplementirani konstrukt odmah prijavi grešku
umesto da prouzrokuje tiho pogrešno ponašanje.
\texttt{KotlinASTBuilder} redefiniše isključivo metode konstrukata obuhvaćeni u implementacije.

\vspace{2px}
Svaki posetilac prima kontekstni objekat \texttt{ctx} koji odražava tačno
ono što je parser pronašao za dato pravilo: ako gramatičko pravilo glasi
\texttt{block : LCURL statements RCURL}, tada \texttt{ctx} ima metode
\texttt{.statements()} i pristupnike za leve i desne vitičaste zagrade.
Obilazak podstabala se vrši rekurzivnim pozivom \texttt{this.Visit(ctx.subRule())}.

\paragraph{Odabir i prosleđivanje na višem nivou.}
Kotlin gramatika grupiše sve konstrukte na nivou datoteke pod pravilo
\texttt{topLevelObject} (funkcije, svojstva, klase, sinoniomi tipova).
Posetilac proverava koji alternativni podizraz nije \texttt{null} i prosleđuje
ga odgovarajućem posetiaocu:

\begin{lstlisting}[language=CSharp, caption={Odabir vrste konstrukta na nivou datoteke}]
public override ASTNode VisitTopLevelObject(
    KotlinParser.TopLevelObjectContext ctx)
{
    if (ctx.functionDeclaration() != null)
        return this.Visit(ctx.functionDeclaration());
    if (ctx.propertyDeclaration() != null)
        return this.Visit(ctx.propertyDeclaration());
    // ... ostale alternative ...
    throw new NotImplementedException(
        $"Unsupported top-level object: {ctx.GetText()}");
}
\end{lstlisting}

\vspace{5px}
\noindent
Isti obrazac primenjuje se i za naredbe unutar bloka - posetilac
\texttt{VisitStatement} razlikuje deklaracije (\texttt{val x = 1}) od izraza
koji se koriste kao naredbe (\texttt{x += 1}, pozivi funkcija).

\paragraph{Prenos greške kao izuzetka.}
Metoda \texttt{Visit} je redefinisana da uhvati \texttt{NullReferenceException}
koji nastaje kada se pristupi opcionalnom gramatičkom elementu koji nije
prisutan u konkretnom ulazu, i pretvori ga u \texttt{SyntaxErrorException}
sa čitljivom porukom. Ovo olakšava dijagnozu problema bez otkrivanja internih
detalja ANTLR obrade.

\newpage
\subsection{Izgradnja čvorova: primer deklaracije}

Sledeći isečak prikazuje kako se gradi LINVAST čvor za deklaraciju
\texttt{val}/\texttt{var} u pet koraka:

\begin{lstlisting}[language=CSharp, caption={Posetilac za deklaraciju svojstva}]
public override ASTNode VisitPropertyDeclaration(
    KotlinParser.PropertyDeclarationContext ctx)
{
    // 1. keyword: "val" or "var"
    string keyword = ctx.VAL() != null ? "val" : "var";

    // 2. variable name
    string name = ctx.variableDeclaration()
                     .simpleIdentifier().GetText();

    // 3. type name (null if no annotation)
    string? typeName = ctx.variableDeclaration()
                          .type()?.GetText();

    // 4. initializer (null if no assignment)
    ExprNode? init = ctx.expression() is not null
        ? this.Visit(ctx.expression()).As<ExprNode>()
        : null;

    // 5. assemble LINVAST nodes
    var specs = typeName is not null
        ? new DeclSpecsNode(ctx.Start.Line, keyword, typeName)
        : new DeclSpecsNode(ctx.Start.Line, keyword);
    var decl = new VarDeclNode(ctx.Start.Line,
                   new IdNode(ctx.Start.Line, name), init);
    return new DeclStatNode(ctx.Start.Line, specs,
               new DeclListNode(ctx.Start.Line, decl));
}
\end{lstlisting}

\noindent
\textbf{Korak 1 : ključna reč.}

\texttt{ctx.VAL()} i \texttt{ctx.VAR()} vraćaju token ako je odgovarajuća
ključna reč prisutna u ulazu, ili \texttt{null} ako nije.
Proverom koja od njih nije \texttt{null} dobija se niska \texttt{"val"} ili
\texttt{"var"} koja se čuva za izgradnju \texttt{DeclSpecsNode}.

\medskip\noindent
\textbf{Korak 2 : ime promenljive.}

Gramatičko pravilo \texttt{variableDeclaration} sadrži \texttt{simpleIdentifier}
koji nosi ime promenljive. Pozivom \texttt{GetText()} dobija se niska, na primer
\texttt{"x"} za deklaraciju \texttt{val x: Int}.

\medskip\noindent
\textbf{Korak 3 : tip (opciono).}

Anotacija tipa (\texttt{: Int}) je opcionalna u Kotlinu i ako je izostavljena,
\texttt{ctx.variableDeclaration().type()} vraća \texttt{null}. Operator
\texttt{?.} osigurava da se u tom slučaju \texttt{typeName} postavi na
\texttt{null} umesto da dođe do izuzetka.

\medskip\noindent
\textbf{Korak 4 : inicijalizator (opciono).}

Ako postoji deo \texttt{= izraz}, rekurzivnim pozivom \texttt{this.Visit}
se izraz prevodi u \texttt{ExprNode}. Koristi se \texttt{is not null} umesto
\texttt{!= null} jer \texttt{ASTNode} preopterećuje operator \texttt{==} pa
obrazac \texttt{is not null} uvek proverava referencu i zaobilazi preopterećeni
operator.

\medskip\noindent
\textbf{Korak 5 : sklapanje čvorova.}

Ekstraktovane vrednosti se slažu odozdo nagore: \texttt{DeclSpecsNode} nosi
ključnu reč i tip, \texttt{VarDeclNode} nosi ime i inicijalizator,
\texttt{DeclListNode} grupiše deklaratore, a \texttt{DeclStatNode} objedinjuje
specifikatore i listu u potpunu naredbu deklaracije.
Ovaj obrazac izdvajanja iz \texttt{ctx} i rekurzivnog sklapanja čvorova
ponavlja se u svakom posetiaocu graditelja.

\subsection{Hijerarhija izraza i prosleđivanje}

Kotlin gramatika definiše prioritet operatora implicitno kroz lanac pravila
od najnižeg ka najvišem prioritetu:

\begin{center}
\texttt{expression}
$\to$ \texttt{disjunction} (\texttt{||})
$\to$ \texttt{conjunction} (\texttt{\&\&})
$\to$ \texttt{equalityComparison} (\texttt{==}, \texttt{!=})
$\to$ \texttt{comparison} (\texttt{<}, \texttt{>}, \ldots)
$\to$ \texttt{additiveExpression} (\texttt{+}, \texttt{-})
$\to$ \texttt{multiplicativeExpression} (\texttt{*}, \texttt{/}, \texttt{\%})
$\to$ \texttt{prefixUnaryExpression} (\texttt{-}, \texttt{!})
$\to$ \texttt{atomicExpression}
$\to$ \texttt{literalConstant} \textbar{} \texttt{simpleIdentifier}
\end{center}

Svaki nivo posetioca prati isti obrazac: ako postoji samo jedan potomak
(nema operatora na tom nivou), rezultat se odmah prosleđuje nižem posetiocu
bez kreiranja novog čvora. Ako postoji više potomaka, gradi se binarni
LINVAST čvor sa operatorom i levim i desnim operandom. Na primer:

\begin{lstlisting}[language=CSharp, caption={Posetilac za aditivni izraz}]
public override ASTNode VisitAdditiveExpression(
    KotlinParser.AdditiveExpressionContext ctx)
{
    if (ctx.multiplicativeExpression().Length == 1)
        return this.Visit(ctx.multiplicativeExpression(0));
    int line = ctx.Start.Line;
    ExprNode left  = this.Visit(ctx.multiplicativeExpression(0))
                         .As<ExprNode>();
    ExprNode right = this.Visit(ctx.multiplicativeExpression(1))
                         .As<ExprNode>();
    var op = ArithmOpNode.FromSymbol(line,
                 ctx.additiveOperator(0).GetText());
    return new ArithmExprNode(line, left, op, right);
}
\end{lstlisting}

\vspace{5px}
\begin{sloppypar}
\noindent
Razmotrimo šta se dešava pri obradi izraza \texttt{3 + 4 * 2}. Parser gradi
CST u kome \texttt{additiveExpression} ima dva potomka tipa
\texttt{multiplicativeExpression}. Pošto levi potomak (\texttt{3}) nema
operator, poziv \texttt{this.Visit} se rekurzivno spušta kroz lanac pravila
sve do \texttt{VisitLiteralConstant}, koji proizvodi \texttt{LitExprNode(3)}.
Desni potomak (\texttt{4 * 2}) ima dva potomka na nivou
\texttt{multiplicativeExpression}, pa \texttt{VisitMultiplicativeExpression}
gradi \texttt{ArithmExprNode("*", LitExprNode(4), LitExprNode(2))} na isti
način. Konačni rezultat je:
\end{sloppypar}

\begin{verbatim}
ArithmExprNode ("+")
+-- LitExprNode (3)
+-- ArithmExprNode ("*")
    +-- LitExprNode (4)
    +-- LitExprNode (2)
\end{verbatim}

\noindent
Prioritet množenja nad sabiranjem je automatski sačuvan strukturom stabla.

\newpage
\subsection{Testiranje}

Ispravnost svake grupe metoda posetioca proverava se kroz NUnit~3 testove
koji su organizovani paralelno sa izvornim datotekama graditelja:
\texttt{DeclarationTests}, \texttt{FunctionTests}, \texttt{StatementTests},
\texttt{ExpressionTests}, \texttt{ClassTests}, \texttt{TypeTests}. Svaki test kreira \texttt{KotlinASTBuilder},
poziva \texttt{BuildFromSource} sa konkretnim izvornim kodom i proverava
strukturu dobijenog stabla. Testovi koriste specijalizovano polazno pravilo
putem preopterećene metode koja prima \texttt{Func<KotlinParser,
ParserRuleContext>}, što omogućava parsiranje samo jednog izraza ili naredbe
bez potrebe za čitavom datotekom:

\begin{lstlisting}[language=CSharp, caption={Primer unit testa za deklaraciju}]
[Test]
public void ValTypedDeclarationTest()
{
    DeclStatNode node = this.GenerateAST("val x: Int")
                            .As<DeclStatNode>();
    Assert.That(node.Specifiers.TypeName, Is.EqualTo("Int"));
    Assert.That(node.DeclaratorList.Declarators.Single().Identifier,
                Is.EqualTo("x"));
}
\end{lstlisting}

\subsection{Deklaracije}
\noindent\textit{Izvorni kod: \texttt{KotlinASTBuilder.Declarations.cs}}

\paragraph{Direktive uvoza.}
\begin{sloppypar}
Gramatičko pravilo \texttt{importHeader} pokriva sve oblike uvoza.
Puna putanja se čita metodom \texttt{ctx.identifier().GetText()} i
čuva u \texttt{ImportNode}. Džoker (\texttt{import kotlin.math.*})
i pseudonim (\texttt{import Foo as Bar}) se parsiraju bez greške,
ali se tiho ignorišu.
\end{sloppypar}

\paragraph{\texttt{val}/\texttt{var} deklaracije.}
Pravilo \texttt{propertyDeclaration} je centralno za sve deklaracije svojstava.
Ime promenljive izdvaja se kroz pod-pravilo \texttt{variableDeclaration}
(\texttt{VisitVariableDeclaration}), koje čita \texttt{simpleIdentifier} i
opcionalnu tipsku anotaciju i vraća \texttt{VarDeclNode}.
Graditelj razlikuje tri moguća slučaja na osnovu prisustva tipske anotacije i
inicijalizatora:

\begin{table}[H]
\centering
\begin{tabularx}{\textwidth}{lllX}
\toprule
\textbf{Ulaz} & \textbf{TypeName} & \textbf{Identifier} & \textbf{Inicijalizator} \\
\midrule
\texttt{val x: Int}      & \texttt{"Int"}    & \texttt{"x"} & \texttt{null} \\
\texttt{val x: Int = 42} & \texttt{"Int"}    & \texttt{"x"} & \texttt{LitExprNode(42)} \\
\texttt{var y = 0}       & \texttt{"var"}    & \texttt{"y"} & \texttt{LitExprNode(0)} \\
\bottomrule
\end{tabularx}
\caption{Ponašanje \texttt{VisitPropertyDeclaration} za različite oblike ulaza}
\end{table}

\noindent
Kada tipska anotacija izostaje (treći red), \texttt{typeName} je \texttt{null} i
\texttt{DeclSpecsNode} se konstruiše samo sa ključnom reči, pa je
\texttt{Specifiers.TypeName} jednako \texttt{"var"}. Ovo je privremeno rešenje jer 
inferencija tipa nije implementirana.

\subsection{Funkcije}
\noindent\textit{Izvorni kod: \texttt{KotlinASTBuilder.Functions.cs}}
\medskip

Pravilo gramatike \texttt{functionDeclaration} pokriva deklaracije funkcija na
nivou datoteke. Graditelj izvlači:

\begin{itemize}
    \item \textbf{ime} - iz \texttt{ctx.identifier().GetText()};
    \item \textbf{povratni tip} - iz \texttt{ctx.type(0).GetText()} ukoliko je
          prisutna dvotačka iza zagrada; u suprotnom se podrazumeva \texttt{"Unit"};
    \item \textbf{parametre} - obilaskom \texttt{ctx.functionValueParameters()},
          što daje \texttt{FuncParamsNode} sa listom \texttt{FuncParamNode}
          objekata; svaki parametar se gradi iz parova (tip, ime);
    \item \textbf{telo} - obilaskom \texttt{ctx.functionBody()} koji proizvodi
          \texttt{BlockStatNode}; ukoliko tela nema (apstraktna ili interfejs
          funkcija), prosledi se \texttt{null}.
\end{itemize}

\begin{sloppypar}
Sintaksa sa izraznim telom (\texttt{fun double(x: Int) = x * 2}) nije
podržana i baca \texttt{NotImplementedException}. Podrazumevane vrednosti
parametara (\texttt{fun greet(name: String = "World")}) se parsiraju bez
greške, ali podrazumevana vrednost se tiho odbacuje.
\end{sloppypar}

\subsection{Deklaracije klasa i interfejsa}
\noindent\textit{Izvorni kod: \texttt{KotlinASTBuilder.Declarations.cs}}
\medskip

Pravilo gramatike \texttt{classDeclaration} pokriva sve oblike deklaracija klasa
i interfejsa. Graditelj proverava prisustvo ključnih reči \texttt{CLASS} i
\texttt{INTERFACE} i na osnovu toga odlučuje koji LINVAST čvor da proizvede:

\begin{lstlisting}[language=CSharp, caption={Posetilac za deklaraciju klase}]
public override ASTNode VisitClassDeclaration(
    KotlinParser.ClassDeclarationContext ctx)
{
    int line = ctx.Start.Line;
    var identifier = new IdNode(line, ctx.simpleIdentifier().GetText());

    TypeNameListNode typeParams = ctx.typeParameters() != null
        ? this.Visit(ctx.typeParameters()).As<TypeNameListNode>()
        : new TypeNameListNode(line);

    TypeNameListNode baseTypes = ctx.delegationSpecifiers() != null
        ? this.Visit(ctx.delegationSpecifiers()).As<TypeNameListNode>()
        : new TypeNameListNode(line);

    var declarations = new List<DeclStatNode>();

    if (ctx.primaryConstructor() != null) {
        FuncParamsNode ctorParams =
            this.Visit(ctx.primaryConstructor()).As<FuncParamsNode>();
        FuncDeclNode ctorDecl = new FuncDeclNode(line, identifier, ctorParams);
        DeclSpecsNode ctorSpecs = new DeclSpecsNode(line, identifier.Identifier);
        declarations.Add(new DeclStatNode(line, ctorSpecs,
                             new DeclListNode(line, ctorDecl)));
    }

    if (ctx.classBody() != null)
        declarations.AddRange(
            this.Visit(ctx.classBody()).As<BlockStatNode>()
                .Children.Cast<DeclStatNode>());

    var declSpecs = new DeclSpecsNode(line, identifier.Identifier);
    var typeDecl  = new TypeDeclNode(line, identifier,
                        typeParams, baseTypes, declarations);

    if (ctx.INTERFACE() != null)
        return new InterfaceNode(line, declSpecs, typeDecl);
    return new ClassNode(line, declSpecs, typeDecl);
}
\end{lstlisting}

\newpage
\paragraph{Primarni konstruktor.}
Kotlin dopušta da se konstruktor definiše direktno u zaglavlju klase:
\newline
\texttt{class Person(val name: String, var age: Int)}. 
\newline
Pravilo gramatike \texttt{primaryConstructor} sadrži \texttt{classParameters} to jest listu parametara
u zagradama. Svaki parametar (\texttt{classParameter}) razlikuje se od parametra
funkcije po tome što može imati prefiks \texttt{val}/\texttt{var} (čime se ujedno
deklariše i svojstvo klase). Primarni konstruktor se preslikava u \texttt{FuncDeclNode}
koji nosi ime klase kao identifikator i parametre kao \texttt{FuncParamsNode}, a
zatim se uvija u \texttt{DeclStatNode} i dodaje na početak liste članova klase.

\paragraph{Nasledni tipovi.}
Lista naslednih tipova iza dvotačke (\texttt{class Dog : Animal}) čita se iz
\texttt{delegationSpecifiers}. Svaki specifikator može biti običan tip
(\texttt{userType}) ili poziv konstruktora (\texttt{constructorInvocation}). U
oba slučaja izvlači se samo ime tipa i čuva u \texttt{TypeNameListNode}.
Eksplicitna delegacija (\texttt{by} izraz) baca \texttt{NotImplementedException}.

\paragraph{Telo klase.}
\texttt{VisitClassBody} obilazi sve \texttt{classMemberDeclaration} i skuplja ih
u \texttt{BlockStatNode}. \texttt{VisitClassMemberDeclaration} prosleđuje
deklaracije funkcija i svojstava postojećim posetiocima
\newline
(\texttt{VisitFunctionDeclaration}, \texttt{VisitPropertyDeclaration}), a za
ugnežđene klase poziva \texttt{VisitClassDeclaration} rekurzivno. Ostali
tipovi članova (sekundarni konstruktori, \texttt{init} blokovi, prateći objekti)
\newline
bacaju \texttt{NotImplementedException}.

\begin{table}[H]
\centering
\begin{tabularx}{\textwidth}{lllX}
\toprule
\textbf{Ulaz} & \textbf{Čvor} & \textbf{Identifier} & \textbf{Napomena} \\
\midrule
\texttt{class Empty} &
    \texttt{ClassNode} & \texttt{Empty} &
    Bez konstruktora ni tela \\
\texttt{class Point(val x: Int)} &
    \texttt{ClassNode} & \texttt{Point} &
    Konstruktor kao prvi član \\
\texttt{class Dog : Animal} &
    \texttt{ClassNode} & \texttt{Dog} &
    \texttt{BaseTypes = ["Animal"]} \\
\texttt{interface Printable \{ \}} &
    \texttt{InterfaceNode} & \texttt{Printable} &
    Prazno telo \\
\bottomrule
\end{tabularx}
\caption{Ponašanje \texttt{VisitClassDeclaration} za različite oblike ulaza}
\end{table}

\subsection{Naredbe kontrole toka}
\noindent\textit{Izvorni kod: \texttt{KotlinASTBuilder.Statements.cs}}
\paragraph{\texttt{if}/\texttt{else}.}
Pravilo \texttt{ifExpression} može imati jednu ili dve grane tela tipa
\texttt{controlStructureBody}. Graditelj proverava prisustvo drugog tela i
gradi binarni ili ternarni \texttt{IfStatNode}:

\begin{itemize}
    \item \texttt{if (x > 0) \{ return x \}} $\to$
          \texttt{IfStatNode(uslov, tada)} (bez \texttt{else});
    \item \texttt{if (x > 0) \{ return x \} else \{ return 0 \}} $\to$
          \texttt{IfStatNode(uslov, tada, inače)}.
\end{itemize}

Telo grane može biti blok \texttt{\{ \ldots \}} (daje \texttt{BlockStatNode}
putem \texttt{VisitBlock}) ili pojedinačan izraz bez vitičastih zagrada (daje
\texttt{ExprStatNode}).

\paragraph{\texttt{while}.}
Pravilo \texttt{whileExpression} uvek daje \texttt{WhileStatNode} sa uslovom
i telom. Ako telo nije prisutno, prosleđuje se prazan \texttt{BlockStatNode}.

\paragraph{\texttt{for}.}
Pravilo \texttt{forExpression} preslikava se u \texttt{ForeachStatNode}.
Graditelj izdvaja promenljivu iteratora iz \texttt{variableDeclaration},
iterable iz \texttt{expression} i telo iz \texttt{controlStructureBody}.
Podržana je i destrukturisana iteracija (\texttt{multiVariableDeclaration}),
gde se više promenljivih pakuje u jedan \texttt{DeclStatNode}.

\paragraph{\texttt{return}.}
Pravilo \texttt{jumpExpression} obuhvata \texttt{return}, \texttt{break},
\texttt{continue} i \texttt{throw}. Implementirani su samo oblici naredbe
\texttt{return}:

\begin{itemize}
    \item \texttt{return 42} $\to$
          \texttt{JumpStatNode(Return, LitExprNode(42))};
    \item \texttt{return} $\to$
          \texttt{JumpStatNode(Return)} bez izraza.
\end{itemize}

\noindent
Naredbe \texttt{break}, \texttt{continue} i \texttt{throw} bacaju
\texttt{NotImplementedException}. Petlja \texttt{do-while} takođe nije
implementirana.

\paragraph{Dodela vrednosti.}
Pravilo \texttt{expression} pokriva i naredbe dodele (\texttt{x = 1},
\texttt{x += 1} itd.). Graditelj proverava prisustvo operatora dodele i
gradi \texttt{ExprStatNode} koji sadrži \texttt{AssignExprNode} sa levim
operandom, operatorom i desnim operandom.

\subsection{Izrazi}
\noindent\textit{Izvorni kod: \texttt{KotlinASTBuilder.Expressions.cs}, \texttt{KotlinASTBuilder.Types.cs}}
\medskip

Svi aritmetički, logički i relacioni izrazi grade se obilaskom hijerarhije
pravila opisane u odeljku~3.3. Podrška po tipu izraza je sledeća:

\begin{itemize}
    \item \textbf{Literali} (\texttt{literalConstant}): celi brojevi, decimalni
          brojevi, logičke vrednosti (\texttt{true}/\texttt{false}) i
          \texttt{null} se preslikavaju u \texttt{LitExprNode} odnosno
          \texttt{NullLitExprNode}. Decimalni brojevi se parsiraju sa
          \texttt{CultureInfo.InvariantCulture} kako bi separator uvek bila
          tačka, bez obzira na lokalna podešavanja operativnog sistema.
    \item \textbf{Identifikatori}: svaki \texttt{simpleIdentifier} daje
          \texttt{IdNode}. Kotlin dozvoljava meke ključne reči (\texttt{get},
          \texttt{set}, \texttt{field} itd.) kao identifikatore, a budući da
          \texttt{VisitSimpleIdentifier} čita tekst tokena putem
          \texttt{ctx.GetText()}, ovi slučajevi se obrađuju automatski.
    \item \textbf{Aritmetika}: \texttt{+}, \texttt{-}, \texttt{*}, \texttt{/},
          \texttt{\%} daju \texttt{ArithmExprNode}.
    \item \textbf{Logika}: \texttt{||}, \texttt{\&\&} daju \texttt{LogicExprNode};
          \texttt{!} (negacija) daje \texttt{UnaryExprNode}.
    \item \textbf{Poređenje}: \texttt{<}, \texttt{>}, \texttt{<=}, \texttt{>=},
          \texttt{==}, \texttt{!=} daju \texttt{RelExprNode}.
    \item \textbf{Zagrade}: \texttt{(izraz)}  zagrade se skidaju a unutrašnji
          izraz se obilazi; u stablu ne postoji zaseban čvor za zagrade.
\end{itemize}

\paragraph{Rezolucija tipova (\texttt{KotlinASTBuilder.Types.cs}).}
Svi tipski nazivi koji se pojavljuju u deklaracijama, parametrima funkcija i
parametrima klasa prolaze kroz četiri posetioca iz ove datoteke:

\begin{itemize}
    \item \texttt{VisitUserType} - spaja segmente kvalifikovanog naziva tačkama
          i proizvodi \texttt{TypeNameNode}, npr.\ \texttt{kotlin.io.File};
    \item \texttt{VisitSimpleUserType} - čita jednostavan naziv tipa i ignoriše
          generičke argumente (\texttt{List<String>} $\to$ \texttt{TypeNameNode("List")});
    \item \texttt{VisitNullableType} - dodaje \texttt{?} na kraj naziva tipa,
          npr.\ \texttt{String?} $\to$ \texttt{TypeNameNode(String?)};
    \item \texttt{VisitTypeReference} - prosleđuje na \texttt{VisitUserType} ili
          rekurzivno skida zagrade; baca izuzetak samo za \texttt{dynamic}.
\end{itemize}

\noindent
Tip funkcije (\texttt{(Int) -> String}) baca \texttt{NotImplementedException}.

\newpage
\subsection{Ilustrativni primer}

Razmotrimo sledeću Kotlin funkciju koja izračunava apsolutnu vrednost celog broja:

\begin{lstlisting}[language=Kotlin, caption={Primer Kotlin funkcije -- apsolutna vrednost}]
fun absolute(x: Int): Int {
    if (x < 0) {
        return -x
    }
    return x
}
\end{lstlisting}

Kada se ovaj isečak prosledi metodi \texttt{BuildFromFile}, \texttt{KotlinASTBuilder}
obilazi stablo parsiranja i gradi sledeću LINVAST reprezentaciju:

\begin{verbatim}
SourceNode
+-- FuncNode (povratni tip: "Int", ime: "absolute")
    +-- FuncParamsNode
    |   +-- FuncParamNode
    |       +-- DeclSpecsNode ("Int")
    |       +-- VarDeclNode ("x")
    +-- BlockStatNode
        +-- IfStatNode
        |   +-- RelExprNode ("<")
        |   |   +-- IdNode ("x")
        |   |   +-- LitExprNode (0)
        |   +-- BlockStatNode
        |       +-- JumpStatNode (Return)
        |           +-- UnaryExprNode ("-")
        |               +-- IdNode ("x")
        +-- JumpStatNode (Return)
            +-- IdNode ("x")
\end{verbatim}

Primetimo da LINVAST stablo ne sadrži informacije o konkretnoj sintaksi Kotlina
(zagrade, dvotačke, ključne reči) - sačuvana je isključivo semantička struktura
koda. Ista analiza primenjena na semantički ekvivalentnu C funkciju proizvela bi
strukturno identično stablo~\cite{ristovic2020}, što demonstrira suštinski cilj
biblioteke.

% -----------------------------------------------------------------------
\section{Budući rad}
\label{sec:future}
% -----------------------------------------------------------------------

Naredni koraci u razvoju Kotlin builder-a obuhvataju kako uklanjanje
postojećih ograničenja, tako i proširenje obuhvata na konstrukte koji
trenutno bacaju \texttt{NotImplementedException}:

\begin{itemize}
    \item \textbf{Lambda izrazi i funkcije višeg reda} - Kotlinovi
          funkcionalni konstrukti zahtevaju proširenje ili aproksimaciju
          LINVAST-a, koji trenutno ne poseduje odgovarajući čvor.
    \item \textbf{Opsezi} - podrška za operator \texttt{..}
          (\texttt{1..10} trenutno baca \texttt{NotImplementedException}).
    \item \textbf{Operatori bezbednosti od null vrednosti} -
          \texttt{?.}, \texttt{?:} (Elvis operator) i \texttt{!!}.
    \item \textbf{Pozivi metoda i pristup poljima} - postfiksne operacije
          poput \texttt{obj.method()} i \texttt{obj.field}.
    \item \textbf{Inferencija tipa} - automatsko određivanje tipa kada
          tipska anotacija izostaje (\texttt{var y = 0} treba da dobije
          \texttt{TypeName = "Int"}, a ne \texttt{"var"}).
    \item \textbf{Sekundarni konstruktori, \texttt{init} blokovi i
          prateći objekti} - konstrukti unutar tela klase koji trenutno
          bacaju izuzetak.
\end{itemize}

\begin{thebibliography}{9}
\bibitem{ristovic2020}
    I.\ Ristović,
    \textit{Jezički invarijantna provera semantičke ekvivalentnosti strukturno
    sličnih segmenata imperativnog koda},
    master rad, Matematički fakultet, Univerzitet u Beogradu, 2020.
\end{thebibliography}

\end{document}
